\section{Simulating the sensor values (the physical model)}\label{sec:6}

BuildSim provides the rooms and no physics, so the air, the people and the plant are all
simulated. The aim is not physical perfection but enough structure (distinct time constants,
daily cycles, genuine responses to actuation) that the controller reacts to meaningful data and
the observer has something real to be wrong about.

\textbf{Every parameter in this section is invented rather than measured, and each is derived from
a stated rule rather than picked.} That distinction matters and is the reason ADR-0008 lists the
derivations: it is what lets Section \ref{sec:13} say precisely which claims the invention invalidates
(the detection-performance numbers) and which it does not (the architectural properties).

\subsection{The six rooms}\label{sec:6.1}

The six rooms are the cluster on corridor A120 of level 0, taken from BuildSim's own
\texttt{level0/\allowbreak floorplan\_data.json}. Every one has a
room → entry → corridor path in BuildSim's walkable graph ending at corridor node \texttt{A120}
and at no other corridor, and the six lie within 59 m of one another along it. Volume is floor area
× an assumed \textbf{3 m ceiling}; the floorplan carries no height. Design flow is the Swedish
workplace guideline, \textbf{7 l/s per person plus 0.35 l/s per m²}, at the room's intended
occupancy.

\textbf{Appendix \ref{sec:C.1}} tabulates all six: area, volume, intended use and schedule, design
occupancy and design flow, with the \textbf{Σ} row the plant is sized from.

Five distinct sizes spanning 18–76.5 m², a 4:1 span in volume, so no two rooms share a time
constant $V/Q$ or a schedule. That matters because R-7's ``unchanged'' is measured against those
differences: the rejected alternative room set (the corridor A170 cluster) had better wall
adjacency but eight of its ten rooms were 69–79 m², and six same-size rooms would pass R-7 more
easily and mean less.

\subsection{The shared air handler, the flow split and the fan laws}\label{sec:6.2}

The rooms do not each get a fan. One air-handling unit \texttt{AHU-A120} serves all six, with one
damper per room. Its flow and power are

\[ Q_\text{fan} = Q_\text{nom}\,s\,(1-\beta), \qquad P = P_\text{nom}\,s^3\,(1+\beta), \]

with speed $s \in [0.2, 1.0]$ and filter blockage $\beta$. The cube in the speed is the fan
affinity law; $\beta$ enters the two with \emph{opposite} signs, which is the whole point: a
blocked filter costs more power for less air, and that asymmetry is what the classifier tests. The
flow is split over the six damper openings in proportion:

\[ Q_i = Q_\text{fan}\,\frac{a_i}{\sum_j a_j}, \qquad a_i = p_i\,A_i, \]

where $p_i$ is the commanded damper position and $A_i$ the damper's authority, taken proportional
to that room's design flow. This is the simplest law under which the dampers genuinely interact,
and the interaction is not incidental: it is the mechanism of Section \ref{sec:7.4}.

Every one of those parameters is invented rather than measured, and each is derived from a stated
rule rather than picked. \textbf{Appendix \ref{sec:C}} gives the full table (parameter, value, derivation,
and the requirement each is first-order on), and §\ref{sec:13.3} says which claims the invention invalidates
and which it does not.

The consequence of the authority rule is worth stating on its own, because it is the design: at
$s=1$ with every damper open each room receives 143 \% of its design flow (the 43 \% of
headroom the table above reserves for compensation), and \emph{no room can take more than its share
of a fixed total}, so \textbf{the only way a room gets more air is if another damper closes.} That is simultaneously what
makes compensation possible and what makes R-7 non-trivial.

\subsection{The air: a well-mixed CO₂ balance per room}\label{sec:6.3}

Each room is one well-mixed volume:

\[ V_i\,\frac{dC_i}{dt} = n_i\,g\cdot 10^{6} - Q_i\,(C_i - C_\text{out}), \]

integrated with forward Euler at a 1 s tick. The $10^{6}$ converts a volumetric generation rate
into the ppm the concentration is measured in: $V$ is m³, $Q$ and $g$ are m³/s, so without it the
production term is a volume fraction while the dilution term is ppm and the two cannot be added.
It is also the factor that makes the sizing come out right: 165.9 l/s holds A145 at roughly
1000 ppm with its 20 design occupants, so the ventilation guideline and the comfort limit agree
about that room, which is the coincidence the fan sizing leans on. There is \textbf{no inter-room
air exchange} and no corridor volume: that was option C's consequence and option C was rejected.

Actuator commands feed back into this equation \emph{through BuildSim}: when the controller opens
a damper, $a_i$ rises, the split gives that room more of $Q_\text{fan}$, $Q_i$ rises, CO₂ falls,
and the next reading reflects it. That is the sense–compute–act loop made concrete.

\subsection{The people}\label{sec:6.4}

Occupancy drives everything else, so it is modelled explicitly as a weekday schedule per room
type: arrivals spread over 07:30–09:00, departures over 16:00–18:00, lunch dips of 30–60 min,
booked meetings (a 0.35 probability per hour for the meeting room, 0.3 per two-hour slot for the
seminar room), 5–8 of 8 desks occupied in the open-plan office, and empty nights and weekends,
with randomness seeded so that no two days are alike but every run is reproducible.

Empty nights do double duty: they are the \textbf{reference condition for the sensor}. With nobody
in the building, CO₂ \emph{must} decay to the outdoor level, so a floor that is wrong overnight is
evidence of a drifting sensor and of nothing else. That is the third of the classifier's
signatures, and it exists because the schedule has nights in it.

\subsection{The faults}\label{sec:6.5}

Four faults, injected \textbf{between the command and its physical effect}, so that the command
record and the physics disagree, which is what gives the monitor something to detect. Four
\emph{mechanisms} against the monitor's three \emph{classifications}: the last two are two ways for
a sensor to be wrong, and the monitor has one word for both of them (§\ref{sec:7.3}).

\begin{longtable}[]{@{}
  >{\raggedright\arraybackslash}p{(\columnwidth - 6\tabcolsep) * \real{0.19}}
  >{\raggedright\arraybackslash}p{(\columnwidth - 6\tabcolsep) * \real{0.13}}
  >{\raggedright\arraybackslash}p{(\columnwidth - 6\tabcolsep) * \real{0.30}}
  >{\raggedright\arraybackslash}p{(\columnwidth - 6\tabcolsep) * \real{0.38}}@{}}
\toprule
\textbf{Fault} & \textbf{Scope} & \textbf{What it does to the physics} & \textbf{Signature it is
meant to leave} \\
\midrule
\endhead
\texttt{STUCK\_DAMPER} & one room & The vane ignores the command and holds its magnitude position
& One residual rises; plant fan power flat\\*
\texttt{BLOCKED\_FILTER} & plant-wide & Scales $Q_\text{fan}$ by $(1-\beta)$ and power by
$(1+\beta)$ & All six residuals rise together; power per unit of flow rises\\*
\texttt{SENSOR\_DRIFT} & one room & Adds a linear, \emph{uncapped} ramp in ppm per hour at the
\emph{sensor}, not in the room & Step-response shape right, overnight floor wrong\\*
\texttt{SENSOR\_OFFSET} & one room & Adds a \emph{fixed} offset in ppm at the sensor from its onset
and holds it there: a recalibration against the wrong reference, or a replacement element
installed without one & The same signature in principle; in practice it confirms within minutes,
long before the overnight floor it would need (§\ref{sec:7.3}) \\
\bottomrule
\end{longtable}

The magnitude means something different for each kind, and the records enforce the range: a frozen
damper position in $[0,1]$; a blockage $\beta$ in the open interval $(0,1)$; a drift rate in ppm per
hour that must be non-zero \emph{and may be negative}; and an offset in ppm under the same rule.
That last permission matters and Section \ref{sec:13.1} returns to it: a negative drift is a fault the
generator can produce and the deployed guarded lower arm does see, because a room reading below
the 420 ppm outdoor concentration betrays itself. The two sensor faults are separate
kinds rather than one parameterised kind because their magnitudes carry different \emph{units} (ppm
per hour against ppm) and one kind holding both would need a fifth field in the spec grammar
below, breaking every scenario already written. Sensor errors on one sensor accumulate; two stuck
positions do not, the last declared one wins.

Faults are declared as \texttt{KIND:TARGET:MAGNITUDE:ONSET}, several separated by \texttt{;} or
newlines, with \texttt{\#} comments ignored, passed to the simulator through \texttt{FAULT\_SPEC}
or \texttt{FAULT\_SPEC\_FILE} so a demo scenario needs no file edit:

\begin{footnotesize}
\begin{verbatim}
STUCK_DAMPER:A115:0.15:T+1800         # vane frozen at the 0.15 floor, 30 min into the run
BLOCKED_FILTER:PLANT:0.4:T+7200       # beta = 0.4; PLANT, because the filter is shared
SENSOR_DRIFT:A124:150:2026-10-05T09:00:00Z   # +150 ppm per hour from that instant
\end{verbatim}
\end{footnotesize}

The ground truth of which faults are active is carried on every simulator snapshot alongside
\emph{both} the true and the measured concentration, so an experiment can always compare what the
detector concluded against what was injected. Sensor noise is Gaussian at
$\sigma = 30$ ppm from a random stream separate from the occupancy schedules, so noise and people
can be varied independently, and that 30 ppm is the same number the CUSUM's $k$ and $h$ are tuned
to, which is why it is not a free parameter.
